\section{Source Code}

\begin{ltgexample}
\Zcref{lst:XML} shows source code written in XML.
\Zcref{line:comment} contains a comment.

\begin{lstlisting}[
  language=XML,
  caption={Example XML Listing},
  label={lst:XML}]
<listing name="example">
  <!-- comment --> (* \label{line:comment} *)
  <content>not interesting</content>
</listing>
\end{lstlisting}
\end{ltgexample}

One can also add \verb+float+ as parameter to have the listing floating.
\Zcref{lst:flXML} shows the floating listing.

\begin{ltgexample}
\begin{lstlisting}[
  % one can adjust spacing here if required
  % aboveskip=2.5\baselineskip,
  % belowskip=-.8\baselineskip,
  float,
  language=XML,
  caption={Example XML listing -- placed as floating figure},
  label={lst:flXML}]
<listing name="example">
  Floating
</listing>
\end{lstlisting}
\end{ltgexample}

One can also typeset JSON as shown in \zcref{lst:json}.

\begin{ltgexample}
\begin{lstlisting}[
  float,
  language=json,
  caption={Example JSON listing -- placed as floating figure},
  label={lst:json}]
{
  key: "value"
}
\end{lstlisting}
\end{ltgexample}

Java is also possible as shown in \zcref{lst:java}.

\begin{ltgexample}
\begin{lstlisting}[
  caption={Example Java listing},
  label=lst:java,
  language=Java,
  float]
public class Hello {
    public static void main (String[] args) {
        System.out.println("Hello World!");
    }
}
\end{lstlisting}
\end{ltgexample}

Listings longer than the remaining space continue on the next page if they are \emph{not} floating.
A floating listing is always placed on a single page.
\Zcref{lst:long} spans multiple pages.

\begin{lstlisting}[
  language=Java,
  caption={Long Java listing spanning multiple pages},
  label={lst:long}]
public class Squares {

    /** Returns the square of 1. */
    public int square1() {
        return 1 * 1;
    }

    /** Returns the square of 2. */
    public int square2() {
        return 2 * 2;
    }

    /** Returns the square of 3. */
    public int square3() {
        return 3 * 3;
    }

    /** Returns the square of 4. */
    public int square4() {
        return 4 * 4;
    }

    /** Returns the square of 5. */
    public int square5() {
        return 5 * 5;
    }

    /** Returns the square of 6. */
    public int square6() {
        return 6 * 6;
    }

    /** Returns the square of 7. */
    public int square7() {
        return 7 * 7;
    }

    /** Returns the square of 8. */
    public int square8() {
        return 8 * 8;
    }

    /** Returns the square of 9. */
    public int square9() {
        return 9 * 9;
    }

    /** Returns the square of 10. */
    public int square10() {
        return 10 * 10;
    }

    /** Returns the square of 11. */
    public int square11() {
        return 11 * 11;
    }

    /** Returns the square of 12. */
    public int square12() {
        return 12 * 12;
    }

    public static void main(String[] args) {
        System.out.println(new Squares().square12());
    }
}
\end{lstlisting}
